\documentclass[10pt,a4paper]{amsart}
\usepackage[margin=19mm]{geometry}
\usepackage{amsmath,amssymb,mathtools}
\usepackage[scr=rsfs,cal=euler]{mathalfa}
\usepackage{xfrac}
\usepackage[unicode,hidelinks,bookmarks=false]{hyperref}
\newcommand{\ie}{i.e.\ }
\newcommand{\cf}{cf.\ }
\newcommand{\resp}{respectively\ }
\newcommand{\Subsection}{\subsection}
\providecommand{\nobreakdash}{\nobreak\hbox{-}}
\setlength{\emergencystretch}{2em}
\multlinegap0pt
\usepackage{bm}
\usepackage{bbm}
\usepackage{dsfont}
\usepackage{xspace}
\usepackage{cancel}
\usepackage{mathtools}
\usepackage{amsthm}
\makeatletter
\@ifundefined{definition}{\newtheorem{definition}{Definition}}{}
\@ifundefined{lemma}{\newtheorem{lemma}[definition]{Lemma}}{}
\@ifundefined{proposition}{\newtheorem{proposition}[definition]{Proposition}}{}
\makeatother
\title[Boundaries in the Instantaneous Formulation of Field Theorie]{Boundaries in the Instantaneous Formulation of Field Theories}
\author{Silvester G.A. Borsboom}
\date{Source: arXiv:2606.31265v1 (CC BY 4.0)}
\begin{document}
\maketitle

\noindent\emph{Source and licence.} Boundaries in the Instantaneous Formulation of Field Theories. Version arXiv:2606.31265v1 (CC BY 4.0), distributed under the Creative Commons Attribution 4.0 International licence, as recorded in the arXiv licence metadata for this exact version and shown on the abstract page https://arxiv.org/abs/2606.31265v1. Full rights notice: licenses/arxiv-2606.31265-CC-BY-4.0.txt.\par\smallskip

\noindent\textit{Editorial scope.} Counted unit: the Proposition whose source label is \texttt{beautifulprop} in the article (source0.tex lines 316--318, source label \texttt{beautifulprop}), delivered together with the article's own complete original proof of it (source0.tex lines 320--368, one \texttt{proof} environment of 49 lines). No mathematics is written, repaired or re-derived: statement and proof are the article's own text line for line. Required context retained uncounted: prop:crucialisom (lines 141--143); kernellemma (lines 263--270); stupidlemma (lines 292--294). Every definition, display and label of those lines is retained unchanged. Omitted from this package and not redistributed: 1--140, 144--262, 271--291, 295--315, 319--319, 369--775. Nothing inside a retained excerpt is omitted or altered: every retained body is delivered exactly as the article's lines are written, so no omission receipt of any kind accompanies this package. Citations: the retained excerpts carry no citation command at all, so no bibliography entry of the article is transcribed and no reference is redistributed. Reference adaptation: none was needed and none was made; every reference inside the delivered unit resolves inside it, so no unresolved reference or citation is produced. Printed slips kept literal: no printed word, symbol, hypothesis or number of the counted statement or of its proof was altered. Layout and class: the article's own class is \texttt{article} with packages including fontenc, graphicx, babel, inputenc, csquotes, cite, endnotes, authblk, geometry, hyperref, dsfont, amsmath, amsfonts, amssymb; the wrapper is the lane's pinned \texttt{amsart} editorial wrapper at 10pt on A4, which restates the theorem-like environments the retained text needs and adds the article's own macro definitions that the retained text needs (none), with extra wrapper packages (bm, bbm, dsfont, xspace, cancel, mathtools). The delivered numbering is the unit's own. Assets: no retained span contains a figure environment, a graphics-inclusion command, a PDF-page inclusion, a table or a TikZ picture; no image file is copied into this package, the package declares no asset dependency and no image file is redistributed with it. Licence: version arXiv:2606.31265v1 of this article is distributed under the Creative Commons Attribution 4.0 International licence, as recorded in the arXiv licence metadata for this exact version and shown on the abstract page https://arxiv.org/abs/2606.31265v1. A full rights notice is delivered at licenses/arxiv-2606.31265-CC-BY-4.0.txt. Characters: every retained excerpt is pure ASCII, so the delivered engine, XeLaTeX with its default Latin Modern text font, reports no missing character.
\begin{proposition}\label{prop:crucialisom}
    Let $Q_t=\Gamma(Y_t)$ denote the instantaneous configuration space at time $t$. Then $\beta_\zeta$ induces an isomorphism $(j^1Q)_t\cong TQ_t$, where $(j^1Q)_t=\{j^1\varphi\circ i_t:\varphi\in\Gamma(Y)\}$ is the collection of holonomic sections of $J^1Y\to M$ restricted to $\Sigma_t$. 
\end{proposition}

\begin{lemma}\label{kernellemma}
For $b_t\in \Gamma(Y_t|_{\partial\Sigma_t})$, we define $Q_t^{b_t}=r_t^{-1}(b_t)$. It then follows that
\begin{align*}
    T_qQ_t^{b_t}
    &= \ker(T_qr_t) = \left\{v\in T_qQ_t :
    v|_{\partial\Sigma_t}=0\right\}.
\end{align*}
\end{lemma}

\begin{lemma}\label{stupidlemma}
    Let $\varphi\in \Gamma(Y)$ be a field history. Then the family of boundary field values $\varphi_t|_{\partial\Sigma_t}=r_t(\varphi_t)$ is constant if and only if $\dot{\varphi}_t|_{\partial\Sigma_t}=0$ for every $t$.
\end{lemma}

\begin{proposition}\label{beautifulprop}
    Let $b$ denote a constant spatial Dirichlet boundary condition, whose restriction to $\partial \Sigma_t$ is denoted $b_t$. Denote by $(j^1Q^b)_t$ the collection of holonomic sections agreeing with $b$ on the spatial boundary, restricted to $\Sigma_t$. Then the isomorphism $(j^1Q)_t\cong TQ_t$ of Proposition \ref{prop:crucialisom}, induced by the jet decomposition map $\beta_\zeta$, restrict to an isomorphism $(j^1Q^b)_t\cong TQ^{b_t}_t$.
\end{proposition}

\begin{proof}
We restrict the jet decomposition map $\beta_\zeta$ to elements of
$(j^1Q^b)_t$. Let $j^1\varphi\circ i_t\in(j^1Q^b)_t$. Since
$\varphi$ agrees with $b$ on the spatial boundary, we have $\varphi_t\in Q_t^{b_t}$.
Moreover, since $b$ is constant, Lemma \ref{stupidlemma} implies $\dot{\varphi}_t|_{\partial\Sigma_t}=0$. It then follows from Lemma \ref{kernellemma} that
\begin{align*}
    \dot{\varphi}_t
    \in
    \ker(T_{\varphi_t}r_t)
    =
    T_{\varphi_t}Q_t^{b_t}.
\end{align*}
Thus $\beta_\zeta(j^1\varphi\circ i_t)
    =(\varphi_t,\dot{\varphi}_t)$ lies in $TQ_t^{b_t}$ under the isomorphism of Proposition \ref{prop:crucialisom}.

Conversely, let $(q,v)\in TQ_t^{b_t}$. Since $Q_t^{b_t}=r_t^{-1}(b_t)$ is a
smooth submanifold of $Q_t$, there exists a smooth curve $\gamma:I\longrightarrow Q_t^{b_t}$ such that $\gamma(0)=q$ and $\dot{\gamma}(0)=v$. For $s\in I$, define a field on the
neighboring slice $\Sigma_{t+s}$ by
\begin{align*}
    \varphi_{t+s}
    :=
    F_s^Y\circ\gamma(s)\circ(F_s^M)^{-1}.
\end{align*}
Clearly $\varphi_t=\gamma(0)=q$. These fields define a field history in a neighborhood of
$\Sigma_t$. Since $\gamma(s)|_{\partial\Sigma_t}=b_t$ and $b$ is
constant with respect to the slicing, we have
\begin{align*}
    \varphi_{t+s}|_{\partial\Sigma_{t+s}}
    =
    F_s^Y\circ b_t\circ(F_s^M)^{-1}
    =
    b_{t+s}.
\end{align*}
Thus $\varphi$ agrees with $b$ on the spatial boundary. Moreover $F_s^*\varphi_{t+s}=\gamma(s)$,
and therefore
\begin{align*}
    \dot{\varphi}_t
    =
    \left.\frac{d}{ds}\right|_{s=0}
    F_s^*\varphi_{t+s}
    =
    \dot{\gamma}(0)
    =
    v.
\end{align*}
Consequently, $(q,v)$ is the image of $j^1\varphi\circ i_t$ under the restriction of
$\beta_\zeta$ to $(j^1Q^b)_t$, using the isomorphism of Proposition
\ref{prop:crucialisom}. By the injectivity of $\beta_\zeta$ this restriction is itself an isomorphism.
\end{proof}

\end{document}
